% 第2章 2.5.2：ℓ^2 上 A(x_n)=(x_n/n) 的像稠密但不闭。
% (a) 目标 y_n=1/n（蓝）与唯一形式逆像 x_n=1（浅色，非平方可和），n=1..10。
% (b) 截断逆像 x^{(N)}=(1,...,1,0,...)：残差 ‖Ax^{(N)}-y‖=(Σ_{n>N}1/n^2)^{1/2}（Python 数值求和）趋于 0，
%     而 ‖x^{(N)}‖=sqrt N 趋于无穷；纵轴为对数刻度，N=1..40。
\begin{tikzpicture}[font=\figurefont\footnotesize]
\pgfplotsset{every axis/.style={width=56mm,height=38mm,axis lines=left,
  x axis line style={draw=FigureStroke,line width=.6pt},y axis line style={draw=FigureStroke,line width=.6pt},
  tick style={draw=FigureStroke},tick label style={font=\figurefont\scriptsize,text=FigureInk},
  label style={font=\figurefont\scriptsize},
  xlabel style={at={(axis description cs:1,0)},anchor=north west,yshift=1mm},
  ymajorgrids,grid style={draw=FigureGrid,line width=.4pt},
  title style={font=\figurefont\footnotesize},
  legend style={draw=none,fill=none,font=\figurefont\scriptsize,cells={anchor=west}}}}
\begin{axis}[name=A,ybar,bar width=1.6mm,xlabel={$n$},xmin=0.3,xmax=10.7,ymin=0,ymax=1.5,
  xtick={1,2,...,10},ytick={0,0.5,1},title={(a) 目标与形式逆像},
  legend style={at={(0.98,0.98)},anchor=north east},legend columns=2]
\addplot[fill=FigureBlue,draw=FigureBlue,bar shift=-0.9mm] coordinates {(1,1) (2,0.5) (3,0.3333) (4,0.25) (5,0.2) (6,0.1667) (7,0.1429) (8,0.125) (9,0.1111) (10,0.1)};
\addplot[fill=FigureYellowFill,draw=FigureYellow,bar shift=0.9mm] coordinates {(1,1) (2,1) (3,1) (4,1) (5,1) (6,1) (7,1) (8,1) (9,1) (10,1)};
\legend{$y_n=1/n$,$x_n=1$}
\end{axis}
\begin{axis}[at={(A.east)},anchor=west,xshift=10mm,xlabel={$N$},xmin=0,xmax=41,ymode=log,ymin=0.1,ymax=10,
  xtick={1,10,20,30,40},ytick={0.1,1,10},yticklabels={$0.1$,$1$,$10$},title={(b) 截断逆像},
  legend style={at={(0.5,-0.3)},anchor=north},legend columns=2]
\addplot[draw=FigureBlue,line width=1.1pt] coordinates {(1,0.8031) (2,0.6284) (3,0.5328) (4,0.4704) (5,0.4258) (6,0.3918) (7,0.3649) (8,0.3428) (9,0.3243) (10,0.3085) (11,0.2948) (12,0.2828) (13,0.2721) (14,0.2626) (15,0.2540) (16,0.2461) (17,0.2390) (18,0.2325) (19,0.2264) (20,0.2208) (21,0.2156) (22,0.2108) (23,0.2063) (24,0.2020) (25,0.1980) (26,0.1942) (27,0.1907) (28,0.1873) (29,0.1841) (30,0.1811) (31,0.1782) (32,0.1754) (33,0.1728) (34,0.1702) (35,0.1678) (36,0.1655) (37,0.1633) (38,0.1612) (39,0.1591) (40,0.1571)};
\addplot[draw=FigureYellow,line width=1.1pt,dash dot] coordinates {(1,1.0000) (2,1.4142) (3,1.7321) (4,2.0000) (5,2.2361) (6,2.4495) (7,2.6458) (8,2.8284) (9,3.0000) (10,3.1623) (11,3.3166) (12,3.4641) (13,3.6056) (14,3.7417) (15,3.8730) (16,4.0000) (17,4.1231) (18,4.2426) (19,4.3589) (20,4.4721) (21,4.5826) (22,4.6904) (23,4.7958) (24,4.8990) (25,5.0000) (26,5.0990) (27,5.1962) (28,5.2915) (29,5.3852) (30,5.4772) (31,5.5678) (32,5.6569) (33,5.7446) (34,5.8310) (35,5.9161) (36,6.0000) (37,6.0828) (38,6.1644) (39,6.2450) (40,6.3246)};
\legend{残差 $\|Ax^{(N)}-y\|$,$\|x^{(N)}\|=\sqrt N$}
\end{axis}
\end{tikzpicture}
